\documentclass[11pt]{letter}
\usepackage[margin=1in]{geometry}
\usepackage{hyperref}

\signature{Elliot Tower}
\address{Independent researcher \\ elliot@elliottower.ai}
\date{July 10, 2026}

\begin{document}
\begin{letter}{Editors \\ \textit{PLOS Computational Biology}}

\opening{Dear Editors,}

I am submitting ``A Pre-Registered Direction Instability Atlas Across Three Perturbation Modalities'' for consideration as a Research Article in \textit{PLOS Computational Biology}.

Whether a perturbation produces consistent effects across cell types determines if laboratory findings transfer to patients, yet no existing metric separates inconsistent effects from merely weak ones. A companion paper under review at this journal (PCOMPBIOL-S-26-02155) introduced direction instability (DI) and validated it on LINCS L1000 drug data. This paper extends DI across three measurement technologies---L1000 transcriptomics, Tahoe-100M single-cell expression, and JUMP Cell Painting morphology---covering ${\sim}$55,000 drugs and genetic perturbations.

Eleven pre-registered hypothesis tests reveal that DI is structured by pharmacological mechanism: drugs targeting core cellular machinery act consistently across cell types, while receptor-signaling drugs produce context-dependent effects. Essential genes produce stereotyped disruption regardless of cell type. A highly selective drug can still have high DI if its target pathway is wired differently across contexts---selectivity alone does not predict transferability. Eight of eleven tests fail their pre-registered criteria, delimiting what the metric captures and preventing overclaiming.

All data, code, pre-registration documents, and results are deposited on Zenodo (DOI: 10.5281/zenodo.21223667). I am the sole author and declare no competing interests. The manuscript is not under consideration at another journal.

\closing{Thank you,}

\end{letter}
\end{document}
